% Options for packages loaded elsewhere
\PassOptionsToPackage{unicode}{hyperref}
\PassOptionsToPackage{hyphens}{url}
\documentclass[
]{article}
\usepackage{xcolor}
\usepackage{amsmath,amssymb}
\setcounter{secnumdepth}{-\maxdimen} % remove section numbering
\usepackage{iftex}
\ifPDFTeX
  \usepackage[T1]{fontenc}
  \usepackage[utf8]{inputenc}
  \usepackage{textcomp} % provide euro and other symbols
\else % if luatex or xetex
  \usepackage{unicode-math} % this also loads fontspec
  \defaultfontfeatures{Scale=MatchLowercase}
  \defaultfontfeatures[\rmfamily]{Ligatures=TeX,Scale=1}
\fi
\usepackage{lmodern}
\ifPDFTeX\else
  % xetex/luatex font selection
\fi
% Use upquote if available, for straight quotes in verbatim environments
\IfFileExists{upquote.sty}{\usepackage{upquote}}{}
\IfFileExists{microtype.sty}{% use microtype if available
  \usepackage[]{microtype}
  \UseMicrotypeSet[protrusion]{basicmath} % disable protrusion for tt fonts
}{}
\makeatletter
\@ifundefined{KOMAClassName}{% if non-KOMA class
  \IfFileExists{parskip.sty}{%
    \usepackage{parskip}
  }{% else
    \setlength{\parindent}{0pt}
    \setlength{\parskip}{6pt plus 2pt minus 1pt}}
}{% if KOMA class
  \KOMAoptions{parskip=half}}
\makeatother
\usepackage{color}
\usepackage{fancyvrb}
\newcommand{\VerbBar}{|}
\newcommand{\VERB}{\Verb[commandchars=\\\{\}]}
\DefineVerbatimEnvironment{Highlighting}{Verbatim}{commandchars=\\\{\}}
% Add ',fontsize=\small' for more characters per line
\newenvironment{Shaded}{}{}
\newcommand{\AlertTok}[1]{\textcolor[rgb]{1.00,0.00,0.00}{\textbf{#1}}}
\newcommand{\AnnotationTok}[1]{\textcolor[rgb]{0.38,0.63,0.69}{\textbf{\textit{#1}}}}
\newcommand{\AttributeTok}[1]{\textcolor[rgb]{0.49,0.56,0.16}{#1}}
\newcommand{\BaseNTok}[1]{\textcolor[rgb]{0.25,0.63,0.44}{#1}}
\newcommand{\BuiltInTok}[1]{\textcolor[rgb]{0.00,0.50,0.00}{#1}}
\newcommand{\CharTok}[1]{\textcolor[rgb]{0.25,0.44,0.63}{#1}}
\newcommand{\CommentTok}[1]{\textcolor[rgb]{0.38,0.63,0.69}{\textit{#1}}}
\newcommand{\CommentVarTok}[1]{\textcolor[rgb]{0.38,0.63,0.69}{\textbf{\textit{#1}}}}
\newcommand{\ConstantTok}[1]{\textcolor[rgb]{0.53,0.00,0.00}{#1}}
\newcommand{\ControlFlowTok}[1]{\textcolor[rgb]{0.00,0.44,0.13}{\textbf{#1}}}
\newcommand{\DataTypeTok}[1]{\textcolor[rgb]{0.56,0.13,0.00}{#1}}
\newcommand{\DecValTok}[1]{\textcolor[rgb]{0.25,0.63,0.44}{#1}}
\newcommand{\DocumentationTok}[1]{\textcolor[rgb]{0.73,0.13,0.13}{\textit{#1}}}
\newcommand{\ErrorTok}[1]{\textcolor[rgb]{1.00,0.00,0.00}{\textbf{#1}}}
\newcommand{\ExtensionTok}[1]{#1}
\newcommand{\FloatTok}[1]{\textcolor[rgb]{0.25,0.63,0.44}{#1}}
\newcommand{\FunctionTok}[1]{\textcolor[rgb]{0.02,0.16,0.49}{#1}}
\newcommand{\ImportTok}[1]{\textcolor[rgb]{0.00,0.50,0.00}{\textbf{#1}}}
\newcommand{\InformationTok}[1]{\textcolor[rgb]{0.38,0.63,0.69}{\textbf{\textit{#1}}}}
\newcommand{\KeywordTok}[1]{\textcolor[rgb]{0.00,0.44,0.13}{\textbf{#1}}}
\newcommand{\NormalTok}[1]{#1}
\newcommand{\OperatorTok}[1]{\textcolor[rgb]{0.40,0.40,0.40}{#1}}
\newcommand{\OtherTok}[1]{\textcolor[rgb]{0.00,0.44,0.13}{#1}}
\newcommand{\PreprocessorTok}[1]{\textcolor[rgb]{0.74,0.48,0.00}{#1}}
\newcommand{\RegionMarkerTok}[1]{#1}
\newcommand{\SpecialCharTok}[1]{\textcolor[rgb]{0.25,0.44,0.63}{#1}}
\newcommand{\SpecialStringTok}[1]{\textcolor[rgb]{0.73,0.40,0.53}{#1}}
\newcommand{\StringTok}[1]{\textcolor[rgb]{0.25,0.44,0.63}{#1}}
\newcommand{\VariableTok}[1]{\textcolor[rgb]{0.10,0.09,0.49}{#1}}
\newcommand{\VerbatimStringTok}[1]{\textcolor[rgb]{0.25,0.44,0.63}{#1}}
\newcommand{\WarningTok}[1]{\textcolor[rgb]{0.38,0.63,0.69}{\textbf{\textit{#1}}}}
\setlength{\emergencystretch}{3em} % prevent overfull lines
\providecommand{\tightlist}{%
  \setlength{\itemsep}{0pt}\setlength{\parskip}{0pt}}
\usepackage{bookmark}
\IfFileExists{xurl.sty}{\usepackage{xurl}}{} % add URL line breaks if available
\urlstyle{same}
\hypersetup{
  hidelinks,
  pdfcreator={LaTeX via pandoc}}

\author{}
\date{}

\begin{document}

\section{The Contract Net Protocol: High-Level Communication and Control
in a Distributed Problem
Solver}\label{the-contract-net-protocol-high-level-communication-and-control-in-a-distributed-problem-solver}

REID G. SMITH, MEMBER, IEEE

Abstract --- The contract net protocol has been developed to specify
problem-solving communication and control for nodes in a distributed
problem solver. Task distribution is affected by a negotiation process,
a discussion carried on between nodes with tasks to be executed and
nodes that may be able to execute those tasks.

We present the specification of the protocol and demonstrate its use in
the solution of a problem in distributed sensing.

The utility of negotiation as an interaction mechanism is discussed. It
can be used to achieve different goals, such as distributing control and
data to avoid bottlenecks and enabling a finer degree of control in
making resource allocation and focus decisions than is possible with
traditional mechanisms.

Index Terms---Artificial Intelligence (AI), connection, cooperation,
distributed problem solving, focus, high-level protocols, negotiation,
resource allocation, task-sharing.

\section{I. INTRODUCTION}\label{i.-introduction}

DISTRIBUTED problem solving is the cooperative solution of problems by a
decentralized and loosely coupled collection of knowledge-sources (KS's)
(procedures, sets of production rules, etc.), located in a number of
distinct processor nodes. The KS's cooperate in the sense that no one of
them has sufficient information to solve the entire problem; mutual
sharing of information is necessary to allow the group, as a whole, to
produce an answer. By decentralized we mean that both control and data
are logically and often geographically distributed; there is neither
global control nor global data storage. Loosely coupled means that
individual KS's spend most of their time in computation rather than
communication.

Such problem solvers offer the promise of speed, reliability,
extensibility, and the potential for increased tolerance to uncertain
data and knowledge, as well as the ability to handle applications with a
natural spatial distribution. There has been much recent interest in
this type of problem solving in the Artificial Intelligence (AI)
community. Its use has been con-

Manuscript received March 4, 1980; revised May 9, 1980. This work was
supported in part by the Advanced Research Projects Agency under
Contract MDA 903-77-C-0322, the National Science Foundation under
Contract MCS 77-02712, and the National Institutes of Health under Grant
RR-00785 on the SUMEX-AIM Computer Facility at Stanford University,
Stanford, CA and by the Defence Research Establishment Atlantic of the
Department of National Defence, Research and Development Branch,
Dartmouth, N.S. Canada.

The author is with the Defence Research Establishment Atlantic,
Dartmouth, N.S., Canada.

sidered in such applications as traffic-light control {[}5{]},
distributed sensing {[}7{]}, and heuristic search {[}10{]}.

In this paper, we present the contract net protocol, a high-level
protocol for communication among the nodes in a distributed problem
solver. It facilitates distributed control of cooperative task execution
(which we call task-sharing {[}9{]}) with efficient internode
communication.

The role of a high-level protocol in a network such as the ARPANET has
been discussed in previous papers (see, for example {[}11{]}).
Traditional communication protocols form a low-level base for
problem-solving communication. They enable reliable and efficient
transmission of bit streams between nodes, but do not consider the
semantics of the information being passed. A high-level protocol assigns
interpretations to the bit streams. It offers a structure that assists
the system designer in deciding what the nodes should say to each other,
rather than how to say it.

We are not primarily concerned with the physical architecture of the
problem solver. It is assumed to be a network of loosely coupled,
asynchronous nodes. Each node contains a number of distinct \(KS's\) .
The nodes are interconnected so that each node can communicate with
every other node by sending messages. No memory is shared by the nodes.
We also assume the existence of a low-level communication protocol to
support reliable and efficient communication of bit streams between
nodes. A functional model of a node is shown in Appendix A.

\section{II. CONNECTION AND CONTRACT
NEGOTIATION}\label{ii.-connection-and-contract-negotiation}

The key issue to be resolved in task-sharing is how tasks are to be
distributed among the processor nodes. There must be a means whereby
nodes with tasks to be executed can find the most appropriate idle nodes
to execute those tasks. We call this the connection problem. Solving the
connection problem is crucial to high performance in a distributed
problem solver. It has two aspects: 1) resource allocation and 2) focus.
Effective resource allocation is achieved by balancing the computational
load among the nodes. It is essential if the maximum speedup possible
from applying multiple nodes to a single overall problem is to be
obtained.

Focus is achieved by effective selection of tasks for allocation to
nodes and by effective selection of KS's for execution of tasks. It is
essential for problems that do not have well-defined

algorithms for their solutions (i.e., problems of the type most often
considered in AI). For such problems, many tasks are typically generated
during the search for solutions; and the execution of many of these
tasks will not lead to a solution. In addition, the most appropriate KS
to invoke for the execution of any given task generally cannot be
identified a priori. The combination of many tasks and many applicable
KS's can lead to a combinatorial explosion. A problem solver must
therefore maintain focus to achieve high performance in practical
applications.

The connection problem can also be viewed from the perspective of an
idle node. It must find another node with an appropriate task that is
available for execution. In our approach, both nodes with tasks to be
executed and nodes ready to execute tasks proceed simultaneously. They
engage each other in discussions that resemble contract negotiation to
solve the connection problem. It is this process that is the basis for
the contract net protocol.

For our purposes, negotiation has four important components: 1) it is a
local process that does not involve centralized control, 2) there is a
two-way exchange of information, 3) each party to the negotiation
evaluates the information from its own perspective, and 4) final
agreement is achieved by mutual selection.

The collection of nodes is referred to as a contract net and the
execution of a task is dealt with as a contract between two nodes. Each
node in the net takes on one of two roles related to the execution of an
individual task: manager or contractor. A manager is responsible for
monitoring the execution of a task and processing the results of its
execution. A contractor is responsible for the actual execution of the
task. Individual nodes are not designated a priori as managers or
contractors; these are only roles, and any node can take on either role
dynamically during the course of problem solving. Typically, a node will
take on both roles, often simultaneously for different contracts. As a
result, nodes are not statically tied to a control hierarchy. This also
leads to more efficient utilization of nodes, as compared, for example,
to schemes that do not allow nodes that have contracted out tasks to
take on other tasks while they are waiting for results.

A contract is established by a process of local mutual selection based
on a two-way transfer of information. In brief, available contractors
evaluate task announcements made by several managers and submit bids on
those for which they are suited. The managers evaluate the bids and
award contracts to the nodes they determine to be most appropriate. The
negotiation process may then recur. A contractor may further partition a
task and award contracts to other nodes. It is then the manager for
those contracts. This leads to the hierarchical control structure that
is typical of task-sharing. Control is distributed because processing
and communication are not focused at particular nodes, but rather every
node is capable of accepting and assigning tasks.

The basic idea of contracting is not new. A rudimentary bidding scheme,
for example, was used for resource allocation in the distributed
computing system (DCS) {[}2{]}, {[}3{]}. We will note the similarities
and differences between that scheme and the contract net protocol as we
proceed. Throughout the paper, reference is made to an experimental
contract net system called CNET. It is a system of INTERLISP {[}12{]}
functions that enables a user to simulate the solution of problems with
a distributed processor.

\section{III. EXAMPLE}\label{iii.-example}

This example is taken from a CNET simulation of a distributed sensing
system (DSS) {[}7{]}. A DSS is a network of sensor and processor nodes
spread throughout a relatively large geographic area. It attempts to
construct and maintain a dynamic map of vehicle traffic in the area.
Construction and maintenance of such a map requires the interpretation
and integration of a large quantity of sensory information received by
the collection of sensor elements.

Use of the contract net protocol in a DSS makes it possible for the
sensor system to be configured dynamically, taking into account such
factors as the number of sensor and processor nodes available, their
locations, and the ease with which communication can be established.

We will examine the negotiation for one particular task, called the
signal task, that arises during the initialization phase of DSS
operation. The task involves gathering of sensed data and extraction of
signal features. The managers for this task are nodes that do not have
sensing capabilities, but do have extensive processing capabilities.
They attempt to find a set of sensor nodes to provide them with signal
features. The sensor nodes, on the other hand, have limited processing
capabilities and attempt to find managers that can further process the
signal features they extract from the raw sensed data.

Recall that we view node interaction as an agreement between a node with
a task to be performed and a node capable of performing that task.
Sometimes the perspective on the ideal character of that agreement
differs depending on the point of view of the participant. For example,
from the perspective of the signal task managers, the best set of
contractors has an adequate spatial distribution about the surrounding
area and an adequate distribution of sensor types. From the point of
view of the signal task contractors, on the other hand, the best
managers are those closest to them in order to minimize potential
communication problems. The ability to express and deal with such
disparate viewpoints is one advantage of the contract net protocol. To
see how the appropriate resolution is accomplished, consider the
messages exchanged between the signal managers and potential signal
contractors. Each signal manager announces its own signal task, using a
message of the sort shown in Fig. 1. Each message in the contract net
protocol has a set of slots for the task-specific information in the
message. The four slots of the task announcement are shown in the
figure. The information that fills the slots is encoded in a simple
language common to all nodes.

The task abstraction is the type of task and the position of the manager
making the announcement. The position enables a potential contractor to
determine the manager to which it should respond. The eligibility
specification indicates that the only nodes that should bid on this task
are those which 1) have sensing capabilities and 2) are located in the
same area as the manager that announced the task. This helps to reduce
extraneous message traffic and bid processing. The bid specifi

To: * indicates a broadcast message.

From: 25

Type: TASK ANNOUNCEMENT

Contract: 22-3-1

Task Abstraction:

TASK TYPE SIGNAL

POSITION LAT 47N LONG 17E

Eligibility Specification:

MUST-HAVE SENSOR

MUST-HAVE POSITION AREA A

Bid Specification:

POSITION LAT LONG

EVERY SENSOR NAME TYPE

Expiration Time:

28 1730Z FEB 1979

Fig. 1. Signal task announcement.

cation indicates the information that a manager needs to select a
suitable set of sensor nodes---the position of the bidder and the name
and type of each of its sensors. Finally, the expiration time is a
deadline for receiving bids.

Each potential contractor listens to the task announcements made by
signal managers. It ranks each announcement relative to the others thus
far received, according to the distance to the manager. Just before the
deadline for the task announcement associated with the perceived nearest
manager, the node submits a bid (Fig. 2). The bid message supplies the
position of the bidder and a description of its sensors. A manager uses
this information to select a set of bidders that covers its area of
responsibility with a suitable variety of sensors, and then awards a
signal contract on this basis (Fig. 3). The award message specifies the
sensors that a contractor must use to provide signal-feature data to its
manager.

\section{IV. THE CONTRACT NET
PROTOCOL}\label{iv.-the-contract-net-protocol}

We now describe the messages of the protocol and the encoding of
information in their slots. We also describe the processing of each
message. The BNF specification of the protocol is presented in Appendix
B. The reader may find it helpful to refer to the specification while
reading the following sections.

\section{A. The Basic Messages}\label{a.-the-basic-messages}

\section{1) Task Announcements}\label{task-announcements}

A node that generates a task normally initiates contract negotiation by
advertising existence of that task to the other nodes with a task
announcement message. It then acts as the manager of the task. A task
announcement can be addressed to all nodes in the net (general
broadcast), to a subset of nodes (limited broadcast), or to a single
node (point-to-point). The latter two modes of addressing, which we call
focused addressing, reduce message processing overhead by allowing
nonaddressed nodes to ignore task announcements after examining only the
addressee slot. The saving is small, but is useful because it allows a
node's communication processor alone to decide whether the rest of the
message should be examined and further processed. It is also useful for
reducing message traffic when the nodes of the problem solver are not
interconnected with broadcast communication channels.

As shown in the example, a task announcement has four main slots. The
eligibility specification is a list of criteria that a node must meet to
be eligible to submit a bid. This slot re-To: 25

From: 42

Type: BID

Contract: 22-3-1

Node Abstraction:

POSITION LAT 62N LONG 9W

SENSOR NAME SI TYPE S

SENSOR NAME S2 TYPE S

SENSOR NAME Tl TYPE T

Fig. 2. Signal task bid.

To: 42

From: 25

Type: AWARD

Contract: 22-3-1

Task Specification:

SENSOR NAME SI

SENSOR NAME S2

Fig. 3. Signal task award.

duces message traffic by pruning nodes whose bids would be clearly
unacceptable. In a sense, it is an extension to the addressee slot.
Focused addressing can be used to restrict the possible respondents only
when the manager knows the names of appropriate nodes. The eligibility
specification slot is used to further restrict the possible respondents
when the manager is not certain of the names of appropriate nodes, but
can write a description of such nodes. \(^{1,2}\)

The task abstraction is a brief description of the task to be executed.
It enables a node to rank the task relative to other announced tasks. An
abstraction is used rather than a complete description in order to
reduce the length of the message. \(^{3}\)

The bid specification is a description of the expected form of a bid. It
enables the manager to specify the kind of information that it considers
important about a node that wants to execute the task. This provides a
common basis for comparison of bids and enables a node to include in a
bid only the information about its capabilities that is relevant to the
task, rather than a complete description. This both simplifies the task
of the manager in evaluating bids and further reduces message traffic.

The expiration time is a deadline for receiving bids. We assume global
synchronization among the nodes. However, time is not critical in the
negotiation process. For example, bids received after the expiration
time of a task announcement are not catastrophic: at worst, they may
result in a suboptimum selection of contractors.

\section{a) The Common Internode
Language}\label{a-the-common-internode-language}

It is useful to encode slot information in a single high-level language
understandable to all nodes We call this a common internode language.
Such a language, along with a high-level programming language (for
transfer of procedures between nodes), forms a common basis for
communicating slot information among the nodes.

\(^{1}\) Note that focused addressing is typically a heuristic process,
since the information upon which it is based may not be exact (e.g., it
may be inferred from prior responses to task announcements).\\
\(^{2}\) We will see that all messages in the protocol that can be
addressed to more than one node have an addressee slot and an
eligibility specification slot to accommodate addressing by name and by
description.\\
\(^{3}\) A numerical priority measure is not, in general, sufficient to
allow potential contractors to rank announced tasks. It assumes first,
that all nodes agree on what constitutes an important task, and second,
that the importance of a task can be captured in a one-dimensional
quantity.

While the contract net protocol offers a framework that specifies the
type of information that is to fill a message slot, it remains the
difficult task of the user to specify the actual content of the slot for
any particular problem domain. In this, sense, the protocol is similar
to AI problem-solving languages like PLANNER {[}4{]}, which supply a
framework for problem solving (e.g., the notions of goal specifications
and theorem patterns), but leave to the user the task of specifying the
content of that framework for any given problem.

CNET does, however, offer the user some additional assistance. It
provides a very simple language, based on an object, attribute, value
representation. The language includes a simple grammar, predefined for
each slot, and a number of predefined domain-independent terms (e.g.,
TASK, TYPE, PROCEDURE, and NAME). The representation, the grammars, and
the domain-independent terms are offered to the user to help him
organize and specify the slot information. He must augment the language
with domain-specific terms (e.g., SENSOR) as needed for the application
at hand.

A message that does not have to be understood by many nodes (e.g.,
messages exchanged by a manager and contractor during execution of a
contract) can be usefully encoded in a private language. This can reduce
both the length of the messages and the overhead required to process
them. In CNET, such ``private'' information is preceded by an ``escape''
character; this allows private information to be inserted in any
message, even one that includes some public information encoded in the
normal manner.

\section{2) Task Announcement
Processing}\label{task-announcement-processing}

In CNET, all tasks are typed. For each type of task, a node maintains a
rank-ordered list of announcements that have been received and have not
yet expired. Each node checks the eligibility specifications of all task
announcements that it receives. This involves ensuring that the
conditions expressed in the specification are met by the node (e.g.,
MUST-HAVE SENSOR). If it is eligible to bid on a task, then the node
ranks that task relative to others under consideration.

Ranking a task announcement is, in general, a task-specific operation.
Many of the operations involved in processing other messages are
similarly task-specific. CNET defines a task template for each type of
task. This template enables a user to specify the procedures required to
process that type of task. In Appendix E we describe the roles of the
required procedures, together with the default actions taken by CNET
when the user chooses to omit a procedure. In the following sections,
whenever reference is made to task-specific actions, the reader may
refer to Appendix E for further details.

\section{3) Bidding}\label{bidding}

This announcement-ranking activity proceeds concurrently with task
processing in a node until the task processor (see Appendix A) completes
processing of its current task and becomes available for processing
another task. At this point, the contract processor is enabled to submit
bids on announced tasks. It checks its list of task announcements and
selects a task on which to submit a bid. If there is only one type of
task, the procedure is straightforward. If, on the other hand, there are
a number of task types available, the node must select one of them. The
current version of CNET selects the most recently received task (older
tasks are more likely to have been already awarded).

An idle node can submit a bid on the most attractive task when either of
the following events occur: 1) the node receives a new task announcement
or 2) the expiration time is reached for any task announcement that the
node has received. At each opportunity, the node makes a (task-specific)
decision whether to submit a bid or wait for further task announcements.
(In the signal task, a potential contractor waits for further
announcements in an attempt to find the closest manager.)

The node abstraction slot of a bid is filled with a brief specification
of the capabilities of the node that are relevant to the announced task.
It is written in the form indicated by the bid specification of the
corresponding task announcement.

The node abstraction slot can also include a number of REQUIRE
statements (e.g., REQUIRE PROCEDURE NAME FFT). Statements of this form
are used by a bidder to indicate that it needs additional information if
it is awarded the task. REQUIRE statements can be made if two conditions
are met: 1) the required objects were not preceded by MUST-HAVE terms in
the eligibility specification of the task announcement and 2) the
objects are transferable; that is, they can be transferred by message.
(A procedure falls into this class, but a hardware device does not.)

The task template is helpful here. If a node receives an announcement
for a type of task with which it is not familiar (i.e., does not have
the template), then it can request the template as a convenient
shorthand for the entire set of procedures associated with that type of
task.

\section{4) Bid Processing}\label{bid-processing}

Contracts are queued locally by the manager that generated them until
they can be awarded. The manager also maintains a rank-ordered list of
bids that have been received for the task. When a bid is received, the
manager ranks the bid relative to others under consideration. If, as a
result, any of the bids are determined to be satisfactory, then the
contract is awarded immediately to the associated bidder. (The
definition of satisfactory is task-specific.) Otherwise, the manager
waits for further bids.

Because a manager is not forced to always wait until the expiration time
before awarding a contract, the average negotiation time for contracts
is reduced over that of DCS {[}2{]}.

If the expiration time is reached and the contract has not yet been
awarded, several actions are possible. The appropriate action is
task-specific, but the possibilities include: awarding the contract to
the most acceptable bidder(s); transmitting another task announcement
(if no bids have been received); or waiting for a time interval before
transmitting another task announcement (if no acceptable bids have been
received). This is in contrast to the traditional view of task
allocation where the most appropriate node available at the time would
be selected.

Successful bidders are informed that they are now contractors for a task
through an announced award message. The task specification slot contains
a specification of the data

needed to begin execution of the task, together with any additional
information requested by the bidder.

\section{5) Contract Processing, Reporting Results, and
Termination}\label{contract-processing-reporting-results-and-termination}

Once a contract has been awarded to a node, it follows the state
transition diagram shown in Appendix C. The data structures shown in
Appendix D form a local context for communication between the contractor
and manager (and other nodes) about the task being performed.

The information message is used for general communication between
manager and contractor during the processing of a contract. (See Section
IV-B.3 for further discussion of this message.)

The report is used by a contractor to inform the manager (and other
report recipients, if any) that a task has been partially executed (an
interim report) or completed (a final report). The result description
slot contains the results of the execution. Final reports are the normal
method of result communication. Interim reports, however, are useful
when generator-style control is desired. A contractor can be set to work
on a task and instructed to issue interim reports whenever the next
result is ready. It then suspends the task until it is instructed by the
manager to continue (with an information message) and produce another
result.

The manager can also terminate contracts with a termination message. The
contractor receiving such a message terminates execution of the contract
indicated in the message and all of its outstanding subcontracts.

\section{6) Negotiation Tradeoffs}\label{negotiation-tradeoffs}

In this section, we discuss choices made in the CNET implementation of
the negotiation process. In the main, our concern has been with problem
solving. We have been more interested in the types of information that
must be passed between nodes than with these aspects of the negotiation
process. As a result, the choices are only tentative and warrant further
detailed analysis.

Because bids are binding and a node is allowed to have more than one bid
outstanding at a time, a node may receive multiple awards. These are
queued for processing in order of receipt. The cost is potentially
slower overall system performance (the load may be less evenly balanced)
than would be the case if multiple awards were prevented.

If nodes could refuse awards (as in DCS {[}3{]}), multiple awards could
be prevented. However, the cost is at least one additional
acknowledgment message per transaction. In some cases, it may be many
additional messages (if an award is refused by several bidders).

Similarly, if a node could only have a single bid outstanding, multiple
awards could be prevented. However, the cost would be significant delay.
Nodes would be forced to remain idle until a task announcement had
expired to find that their bids had been rejected, and have to start the
process again. This could lower overall system performance.

The above delay could be reduced in some instances by explicitly
informing unsuccessful bidders that their bids had been refused. The
cost, however, would likely be a very large increase in message traffic,
assuming that there are several bidders per task. In addition, it would
only lower the delay for contracts that were awarded before the
expiration times of their task announcements were reached.

We have allowed a node to bid at intervals related to the receipt of
task announcements rather than at fixed intervals. It is intuitively
appealing, offers a reasonable compromise between message traffic and
delay in allocating tasks, and has exhibited good performance in
experience to date with CNET.

We have chosen to allow a node to submit, at most, a single bid at each
opportunity. This reduces message traffic and the possibility that a
single node could bid on far more tasks than it could process. For the
same reasons, we allow only idle nodes to submit bids. Because of these
choices, contracts are not-centrally notarized as they were in DCS
{[}2{]}. This further reduces message traffic and maintains the
distributed nature of the negotiation process.

The current version of CNET uses a nonpreemptive scheduler in each node.
It would appear to be useful, however; to allow preemptive scheduling
instead of simply queueing contracts in order of receipt. This would
help avoid the situation where time-critical tasks are not executed soon
enough because less important tasks are queued ahead of them. The
difficult question here is determining the criteria for preemption and
providing a place for them in the common internode language. We are
currently exploring this issue.

\section{B. Complications and
Extensions}\label{b.-complications-and-extensions}

In the following sections we consider some problems with the basic
negotiation mechanism and present extensions to solve these problems.
The extensions are evolving as we gain more experience with using the
protocol in practical applications.

\section{1) Immediate Response Bids}\label{immediate-response-bids}

We have, thus far, discussed the negotiation process under the
assumption that a node cannot submit a bid until it goes idle and is
actually ready to process a new contract. This strategy can, however,
lead to difficulty. For instance, a node that issues a task announcement
may not receive any bids for one of several reasons: 1) there are no
idle nodes; 2) some node is both idle and eligible, but ranks the task
too low; or 3) no node is capable of working on the task even if it were
idle (as may happen if the eligibility specification is too stringent or
no node has the data necessary for executing the task). In the first two
cases, the task announcement may be usefully reissued until a bid is
obtained from an idle node; in the third case it would be pointless.
Therefore, a node requires a way of determining what caused the lack of
response.

A class of bids we call immediate response bids offers a mechanism for
doing this. Three such bids have been identified, allowing a node to
indicate that it is eligible but BUSY, that it is INELIGIBLE, or that it
gave the task a LOW RANKING. A manager can specify that nodes are to
respond in any of these cases or to respond in a subset of the cases
(e.g., respond if eligible, but busy).

A node receiving a task announcement whose.bid specification asks for an
immediate response bid does not deal with

the announcement in the usual way (ranking it immediately, but waiting
until it is idle to submit a bid), but instead responds immediately with
either a standard bid or the appropriate special form.

The immediate response mechanism permits a manager to take a more
appropriate course of action if a task announcement elicits no bids. The
normal procedure is to simply reissue the task announcement. If this
continues to elicit no bids, then the manager can specify an immediate
response bid. If the response is uniformly BUSY, then the manager can
wait and reissue the task announcement later. If all nodes are
INELIGIBLE, then the manager may loosen the eligibility specification.
If all nodes gave the task a LOW RANKING, the manager can wait and
reissue, in the hope that the more interesting tasks currently occupying
the nodes will soon be done. The manager may also decide to execute the
task itself when it becomes idle.

\section{2) Directed Contracts}\label{directed-contracts}

The normal contract negotiation process can be simplified in some
instances, with a resulting enhancement in the efficiency of the
protocol. If a manager knows exactly which node is appropriate for
execution of a task, a directed contract can be awarded. No task
announcement is made and no bids are submitted. Instead, a directed
award message is sent to the selected node without negotiation.

In addition to the low-level acknowledgment used by an underlying
communication protocol to ensure reliable communication, the contract
net protocol uses a further high-level acknowledgment message for
directed award messages. This allows a refusal or negative
acknowledgment of the form, I can't execute this contract because . . .,
in addition to the low-level negative acknowledgment of the form, I
didn't receive your message correctly.

If the node that receives a directed award either does not meet the
eligibility specification or does not have a template for the task, then
it transmits a refusal to the manager. Otherwise, the node uses
task-specific criteria to determine whether to transmit an acceptance or
refusal to the manager. The refusal message contains a refusal
justification slot that specifies the reasons for the refusal (e.g.,
INELIGIBLE or NO TEMPLATE).

The action taken by a manager upon receipt of a refusal is also
task-specific, in general.

\section{3) Request and Information
Messages}\label{request-and-information-messages}

If a simple transfer of information is required, then a request-response
sequence can be used without further embellishment. The request message
is used to encode straightforward requests for information. The
information message is used to respond to a request, and for
manager/contractor communication during the execution of a contract
(Section IV-A.5). This message is also used in result-sharing {[}9{]}, a
style of distributed problem solving in which nodes cooperate by sharing
partial results.

If a node that receives a request message meets the eligibility
specification and has the necessary information, then it responds with
an information message. No task-specific decisions need to be made. If a
node that receives an information message meets the eligibility
specification, then the action it takes is task-specific. If the
contract named in the message is one that the node is processing, then
the appropriate action is determined by a procedure associated with that
contract. Otherwise, the node takes a default action. (The latter
possibility arises in result-sharing because nodes communicate partial
results to other nodes without being requested to do so, and without
being explicitly linked by a contract.)

\section{4) Node Available Messages}\label{node-available-messages}

The protocol has been designed to allow a reversal of the normal
negotiation process. When the processing load on a net is high, most
task announcements will not be answered with bids because all nodes will
be busy. Hence, the protocol includes a node available message. A node
that transmits such a message is idle and searching for a task to
execute. In this case, the eligibility specification is a list of
criteria that the node will look for in a task. The node abstraction is
a brief specification of the capabilities of the node and the expiration
time is a deadline for receiving an award.

When a node receives a node available message, it tries to determine if
it can match the node with a task that it has announced, but not yet
awarded. First, it tries to match the eligibility specification in the
node available message against the task abstractions of the tasks that
it has waiting. A task for which this match succeeds is of interest to
the volunteering node. The manager then tries to match the eligibility
specification of the task against the node abstraction of the node
available message. If this match succeeds, then the node is the sort
which is required for the task. Hence, even in this case we have the
mutual selection aspect, as the volunteering node selects a task and the
task specifies the criteria necessary in a volunteer. If the two matches
are successful, the manager sends a directed award message to the
volunteer. If a node has several tasks for which both matches succeed,
then CNET selects the oldest task (it is least likely to be awarded in
the normal way).

In CNET, a node maintains a list of unexpired node available messages.
Before it announces a new task, a node first tries to find a suitable
volunteer in the list. If several volunteers are suitable, the newest
volunteer is selected (it is least likely to have received an award).

A node can thus acquire a contract in one of two ways: it can wait for a
suitable task announcement and submit a bid; or it can transmit a node
available message and wait for a directed award. The decision as to
which method to use is net-load dependent. If the net is not heavily
loaded, then the use of the task announcement is warranted in all cases,
since the availability of a task to be executed is the event of primary
importance. If the net is heavily loaded, then unlimited use of task
announcements serves only to saturate the available communication
channels. In this case, the availability of an idle node is the event of
primary importance and the use of node available messages is preferred.

In CNET, a node selects one of the two modes based on experience with
its own task announcements. The scheme is rudimentary, however, and
warrants further analysis.

\section{V. DYNAMIC DISTRIBUTION OF
INFORMATION}\label{v.-dynamic-distribution-of-information}

The contract net protocol enables dynamic distribution of information
(i.e., procedures and data) via three methods. First, a node can
transmit a request directly to another node for the transfer of the
required information. The response is the information requested (e.g.,
the code for a procedure). Second, a node can broadcast a task
announcement in which the task is a transfer of information. A bid on
the task indicates that the bidder has the information and is willing to
transmit it. Finally, a node can note, in its bid on a task, that it
requires particular information in order to execute the task. The
manager can then send the required information in the award message if
the bid is accepted.

Dynamic distribution enables effective use of available computational
resources: a node that is standing idle because it lacks information
required to execute a previously announced task can acquire that
information as indicated above. This also means that nodes do not have
to be preloaded with extensive amounts of information which may or may
not prove useful. Dynamic distribution also facilitates the addition of
a new node to an existing net; the node can dynamically acquire the
procedures and data necessary to allow it to participate in the
operation of the net. This is especially useful in the distributed
sensing application.

Dynamic distribution works, first, because all nodes know the syntax of
the contract net protocol, which enables them to identify the slots that
contain task-dependent information (e.g., the eligibility specification
of a task announcement); second, because they know how to parse the
common internode language, which enables them to identify the
information in the slots that they do not possess; and third, because a
negotiation process is used to effect task distribution, which makes it
possible for a new node to begin participating in the operation of the
net by listening to the messages being exchanged (specifically, task
announcements) and submitting bids. This may be contrasted with more
traditional task allocation schemes that explicitly assign tasks to
nodes. In such schemes, new nodes must be explicitly linked to other
nodes in the network.

\section{VI. DISCUSSION}\label{vi.-discussion}

We have presented a high-level protocol for distributed problem solving.
The protocol facilitates distributed control of cooperative task
execution with efficient internode communication. It has been designed
to provide a more powerful mechanism for connection than is available in
current problem-solving systems. The message types have been selected to
capture the types of interactions that arise in a task-sharing approach
to distributed problem solving. The message slots have been selected to
capture the types of information that must be passed between nodes to
make these interactions effective.

The connection of nodes with tasks to be executed and nodes that can
execute those tasks is more general than load balancing. The difference
can be illustrated by comparing the information placed in the slots of
the task announcement with that used in DCS. In that system, a task
announcement (request for quotation) contains only the name of the
process to be executed and the amount of memory required. A bid is a
measure of the excess memory available in a node. Note, however, that
this is only one of the possible dimensions along which tasks and nodes
can be evaluated. In general, such evaluation is task-dependent (as we
saw in Section III) and may require a different form of information for
each task. Recognizing this, the contract net protocol takes a wider
perspective and allows a range of descriptions to be passed between the
manager and bidders. In more general terms, since our concern is with
problem solving, more attention is paid to the type of information that
must be transferred between nodes to solve the connection problem. The
connection that is effected with the contract net protocol is an
extension to the pattern-directed invocation used in many AI programming
languages (see {[}8{]} for a more in-depth discussion).

There are several reasons for adopting a distributed approach to problem
solving. These include speed, reliability, extensibility, and the
ability to handle applications that have a natural spatial distribution.
The design of the contract net protocol attempts to ensure that these
potential benefits are indeed realized.

In order to achieve high speed we wish to avoid bottlenecks. Such
bottlenecks can arise in two primary ways: by concentrating
disproportionate amounts of computation or communication at central
resources and by saturating available communication channels so that
nodes must remain idle while messages are transmitted.

To avoid bottlenecks we distribute control, data, and \(KS'\) s. Control
is distributed through the use of negotiation to achieve connections for
task distribution. Every node is capable of assigning and accepting
tasks, and managers and contractors simultaneously seek each other out.
Data and \(KS'\) s are also distributed dynamically as part of the
negotiation process or with a request-response mechanism.

The contract net protocol includes several elements aimed at avoiding
communication channel saturation. First, the information in task
announcements (like eligibility specifications) helps eliminate
extraneous message traffic. Similarly, bid messages can be kept short
and ``to the point'' through the use of the bid specification mechanism.
Second, specialized interactions, like directed contracts and requests,
reduce communication for transactions that do not require the complexity
of negotiation. Third, the protocol enables dynamic distribution of
information on an as-required basis. Finally, it provides a very general
form of guidance in determining appropriate partitioning of problems:
the notion of tasks executed under contracts suggests relatively large
task sizes. This is important if the speedup obtained from distribution
is not to be outweighed by the effort required to effect that
distribution.

Distribution of control also enhances reliability and permits graceful
degradation of performance in the case of individual node failures.
There are no nodes whose failure can completely block the contract
negotiation process. In addition, recovery from the failure of a node is
aided by the presence of explicit links between managers and their
contractors. The failure of any contractor, for example, can be detected
by its manager;

the contract for which it was responsible can then be reannounced and
awarded to another node.

The use of a negotiation process for allocation of nodes to tasks
facilitates the extension of an existing net to include new nodes. Such
nodes can begin to participate in the operation of the net without first
explicitly informing the other nodes of their presence. New node entry
is also facilitated by the common internode language and dynamic
distribution of information.

The ability to handle applications with a natural spatial distribution
is facilitated by the local nature of negotiation. A net is able to
configure itself dynamically according to the positions of the nodes and
the ease with which they can establish communication.

The protocol is best suited to problems in which it is appropriate to
define a hierarchy of tasks (e.g., heuristic search) and a hierarchy of
levels of data abstraction (e.g., audio or video signal interpretation).
Such problems lend themselves to decomposition into a set of relatively
independent subtasks with little need for global information or
synchronization. Individual subtasks can be assigned to separate
processor nodes; these nodes can then execute the subtasks with little
need for communication with other nodes.

\section{VII. CONCLUSION}\label{vii.-conclusion}

The main contribution of the contract net protocol is the mechanism it
offers for structuring high-level interactions between nodes for
cooperative task execution. It stresses the utility of negotiation as an
interaction mechanism. Negotiation can be used at different levels of
complexity. At one extreme, it is a means of achieving task distribution
with distributed control and shared responsibility for tasks to maintain
reliability and avoid bottlenecks. At the other extreme, the two-way
transfer of information and mutual selection attributes of negotiation
make possible a finer degree of control in making resource allocation
and focus decisions than is possible with traditional mechanisms.

\section{APPENDIX A}\label{appendix-a}

\section{CONTRACT NODE ARCHITECTURE}\label{contract-node-architecture}

We view a node as having four major functional components (Fig. 4): a
local database, a communication processor that handles low-level message
traffic with other nodes, a task processor that carries out the
computation associated with user tasks, and a contract processor that
processes high-level protocol messages and manages node resources. It is
assumed that the functions of the three processors are carried out
concurrently.

\section{APPENDIX B}\label{appendix-b}

\section{CONTRACT NET PROTOCOL
SPECIFICATION}\label{contract-net-protocol-specification}

The BNF specification of the contract net protocol is shown in Fig. 5.
Nonterminal symbols are enclosed by ``( ),'' terminal symbols are
written without delimiters, and nonterminals whose specific expansion is
not germane to the discussion are enclosed by ``{[} {]}.'' Slots that
are to be filled with information encoded in the common internode
language are enclosed in ``\{ \}.'' Message types that need not be
included in a basic implementation are followed by ``*.''

\textit{[Image omitted: images/07ff6467479ef95356ac26162fd76a9a3d5d554f5a87e875c4fb8ba5a4336b61.jpg]}

flowchart

\begin{Shaded}
\begin{Highlighting}[]
\NormalTok{graph TD}
\NormalTok{    A["CONTRACT NODE"] {-}{-}\textgreater{} B["COMMUNICATION PROCESSOR"]}
\NormalTok{    B {-}{-}\textgreater{} C["LOCAL DATA BASE"]}
\NormalTok{    C {-}{-}\textgreater{} D["CONTRACT PROCESSOR"]}
\NormalTok{    C {-}{-}\textgreater{} E["TASK PROCESSOR"]}
\end{Highlighting}
\end{Shaded}

Fig. 4. Contract node architecture.

⇒

⇒ {[}line-header{]} {[}identifier{]} {[}time{]} {[}acknowledge{]} ⇒
{[}error-control{]} {[}line-trailer{]} ⇒ {[}net-address{]} \textbar{}
{[}subnet-address{]} \textbar{} {[}node-address{]} ⇒ {[}node-address{]}
⇒ {[}cctext{]} \textbar{} ⇒ \textbar{} \textbar{} \textbar{} \textbar{}
\textbar{} \textbar{} \textbar{} \textbar{} \textbar{} ⇒
TASK-ANNOUNCEMENT {[}name{]} \{task-abstraction\}
\{eligibility-specification\} \{bid-specification\}
{[}expiration-time{]} ⇒ BID {[}name{]} \{node-abstraction\} ⇒
ANNOUNCED-AWARD {[}name{]} \{task-specification\} ⇒ DIRECTED-AWARD
{[}name{]} \{task-abstraction\}* \{eligibility-specification\}
\{task-specification\} ⇒ ACCEPTANCE {[}name{]}\emph{ ⇒ REFUSAL
{[}name{]} \{refusal-justification\} ⇒ INTERIM-REPORT {[}name{]}
\{result-description\} FINAL-REPORT {[}name{]} \{result-description\} ⇒
TERMINATION {[}name{]}} ⇒ NODE-AVAILABLE \{eligibility-specification\}*
\{node-abstraction\} {[}expiration-time{]} ⇒ REQUEST {[}name{]}
\{eligibility-specification\}* \{request-specification\} ⇒ INFORMATION
{[}name{]} \{eligibility-specification\}*
\{information-specification\}\\
Fig. 5. Contract net protocol specification.

One possible set of low-level message slots is shown for completeness.
Other variations may be dictated by the specific communication
architecture for which the protocol is implemented. Briefly: line-header
delimits the beginning of a message; identifier is a unique identifier
for the message; time specifies the time at which the message is
transmitted; acknowledge is set only if acknowledgment of receipt of the
message is required (this slot is included because delivery of some
messages is not essential to the operation of the protocol (e.g.,
broadcast task announcements) and acknowledgments for such messages
might greatly increase message traffic); error-control is used by an
addressee to determine that the message has been correctly received; and
line-trailer delimits the end of the message. Three forms of addressee
slot are shown. These are used for general broadcast, limited broadcast,
and point-to-point messages, respectively.

There are two forms of text slot. cctext (communication control text)
indicates that the message is for checking net integrity, acknowledging
receipt of messages, maintaining

basic net communication data, such as routing tables, and so on.
Messages of the form Hello and I heard you, exchanged periodically by
IMP's in the ARPANET as status checks, fall into this category.

pstext (problem-solving text) indicates that the message is for
high-level problem solving (see Section IV).

\section{APPENDIX C}\label{appendix-c}

\section{CONTRACT PROCESSING STATES}\label{contract-processing-states}

Contracts are processed according to the state transition diagram shown
in Fig. 6, which uses the terminology of operating systems {[}1{]}. The
contract processor of a node (Appendix A) is responsible for moving a
contract through the states of the diagram. The standard progression
through the states is shown by solid lines. Optional progressions are
shown by dashed lines.

When a contract is awarded to a node (IN), it is placed in the READY
state where it waits until the task processor is available. At that
time, the contract is placed in the EXECUTING state and its processing
is started. If subcontracts are generated, they are placed in the
ANNOUNCED state. A subcontract may either be awarded to another node
(OUT) or to this node (READY).

If execution cannot continue on the contract until some event occurs
(e.g., receipt of a subcontract report), then the contract is placed in
the SUSPENDED state; and another contract in the READY state is selected
for processing. When the awaited event occurs, the contract is
transferred back to the READY state where it awaits further processing.

When processing of the contract has been completed, it is placed in the
TERMINATED state; and another contract in the READY state is selected
for processing. Recently completed contracts are held in the TERMINATED
state in order to facilitate recovery from the failure of a manager.

If a contract in the EXECUTING state is moved to either the SUSPENDED or
TERMINATED state and if there are no contracts in the READY state, then
the node attempts to acquire a new contract---either by making a bid on
a recent task announcement or by transmitting a node available message
(Section IV).

The scheduler used by a contract processor in the current version of
CNET is nonpreemptive and gives priority to a contract whose execution
can continue (e.g., due to receipt of a report) over a contract whose
execution has not yet been started.

\section{APPENDIX D}\label{appendix-d}

\section{CONTRACT SPECIFICATION}\label{contract-specification}

A node maintains a structure of the form shown in Fig. 7 for each task
for which it is the contractor. Such a structure forms a local context
for the execution of a task.

name is a unique identifier for the contract. manager is the node that
generated the task. report recipients are the nodes to which reports for
the contract are to be sent. The default report recipient is the
manager. related contractors are the nodes that are working on related
contracts (e.g., subcontracts of the same contract).

The report recipients and related contractors slots facilitate more
flexible communication and cooperation. They give IN

\textit{[Image omitted: images/a3e54c3679d294c1eeb6aec16fe211c1a471fcd82b90a85a0e917743454324b8.jpg]}

flowchart

\begin{Shaded}
\begin{Highlighting}[]
\NormalTok{graph TD}
\NormalTok{    IN {-}{-}\textgreater{}|IN| A["RECEIVE AWARD"]}
\NormalTok{    A {-}{-}\textgreater{} B["WAIT FOR TASK PROCESSOR"]}
\NormalTok{    B {-}{-}\textgreater{} C["OBTAIN TASK PROCESSOR"]}
\NormalTok{    C {-}{-}\textgreater{} D["EXECUTING"]}
\NormalTok{    D {-}{-}\textgreater{} E["WAIT FOR REPORT"]}
\NormalTok{    E {-}{-}\textgreater{} F["COMPLETE CONTRACT"]}
\NormalTok{    F {-}{-}\textgreater{} G["MAKE AWARD"]}
\NormalTok{    G {-}{-}\textgreater{} OUT}
\NormalTok{    D {-}{-}\textgreater{} H["ANNOUNCED TERMINATED SUSPENDED"]}
\NormalTok{    H {-}{-}\textgreater{} I["GENERATE SUBCONTRACT"]}
\NormalTok{    I {-}{-}\textgreater{} J["RELEASED"]}
\NormalTok{    J {-}{-}\textgreater{} A}
\end{Highlighting}
\end{Shaded}

Fig. 6. Contract processing states.

\textit{[Image omitted: images/e9569e61ab53ba526643cd810abfc9da858d0cd5c1da8c950e2d3ee3dad98256.jpg]}

text\_image

⇒ {[}name{]} {[}manager{]} {[}report-recipients{]}
{[}related-contractors{]}

{[}results{]}

Fig. 7. Contract structure.

node an explicit indication of other nodes in the net with which it may
be useful to communicate (e.g., to enable use of focused addressing in
task announcements) and help to reduce message traffic. Messages from
one related contractor to another, for example, can be addressed
directly and therefore do not have to follow the communication paths of
the control hierarchy.

task is filled with a structure of the form shown in Fig. 8. The
structure lists the type of the task, its specification, and the
procedures that are required to handle the task in a contract net. These
procedures are described in Appendix E.

results are the outcome of executing a task. They are transmitted to the
report recipients of the contract. subcontract-list is the collection of
subtasks generated from the initial contract. This slot is filled with a
list of structures of the form shown in Fig. 9. Such structures contain
the information held by a manager for a subtask.

name is a unique identifier for the subcontract. contractor is the name
of the node responsible for its processing. task is the task structure
(Fig. 8). results are the outcome of executing the subtask. predecessors
are the names of subcontracts that are preconditions for the subcontract
and successors are the names of subcontracts for which the subcontract
is a precondition. Predecessors and successors are necessary for tasks
that must be performed in a particular order, such as those that occur
in planning applications {[}6{]}. A subcontract will not be announced
(Section IV-A.1) until its predecessors have been completed.

\section{APPENDIX E}\label{appendix-e}

\section{TASK PROCESSING PROCEDURES}\label{task-processing-procedures}

In this appendix, we discuss the task-specific procedures, the roles
they play, and the default actions taken in CNET when no procedure is
provided by the user. In the terms of the functional model of a node
presented in Appendix A, the execution procedure runs in the task
processor. The rest of the procedures run in the contract processor.

Announcement Procedure: Called each time a subtask is generated. Its
role is to determine whether the task should be announced or awarded
directly. In either case, it must determine the addressee. It must also
compute the information re-

⇒ {[}name{]} {[}type{]} {[}specification{]} {[}announcement-procedure{]}
{[}announcement-ranking-procedure{]} {[}bid-procedure{]}
{[}bid-ranking-procedure{]} {[}award-procedure{]}
{[}acknowledgment-procedure{]} {[}refusal-processing-procedure{]}
{[}report-acceptance-procedure{]} {[}termination-procedure{]}
{[}information-acceptance-procedure{]} {[}execution-procedure{]} Fig. 8.
Task structure.

=\textgreater{[}name{]} {[}contractor{]}

{[}results{]} {[}predecessors{]} {[}successors{]}\\
Fig. 9. Subcontract structure.

quired to fill the slots of a task announcement or directed award and
indicate the names of any predecessor subcontracts. The default is to
announce the task to all nodes with a task abstraction that states the
task type and a default expiration time.

Announcement Ranking Procedure: Called when a task announcement is
received. Its role is to rank the new announcement relative to others
under consideration. The default is to give the most recently received
task announcement the highest ranking.

Bid Procedure: Called when a node has an opportunity to submit a bid.
Its role is to decide whether or not to submit a bid and to add any
common internode language statements to those called for in the bid
specification of the associated task announcement. No default action is
taken.

Bid Ranking Procedure: Called when a bid is received. Its role is to
rank the new bid relative to others under consideration and determine
whether the contract should be awarded to any of the current bidders.
The default is to insert the new bid into the list of bids so that bids
that REQUIRE the least information (Section IV-A.3) are ranked highest.

Award Procedure: Called when the expiration time for a task announcement
is reached and the contract has not yet been awarded. Its role is to
decide what action is to be taken. The default is to award the contract
to the bidder that REQUIRE's the least additional information (or the
first bidder in case of a tie) or to transmit another task announcement
if no bids have been received.

Acknowledgment Procedure: Called when a node receives a directed
contract for which it is eligible and has a task template. Its role is
to decide whether to accept or reject the contract. The default is to
accept.

Refusal Processing Procedure: Called when a refusal message is received.
Its role is to determine what action is to be taken. The default is to
transmit a task announcement.

Report Acceptance Procedure: Called when a report message is received.
Its role is to determine what to do with the contents of the result
description slot. The default is to add the contents of the slot to the
results slot of the contract.

Termination Procedure: Called when a contract that has been terminated
is about to be deleted from the local database.

(CNET maintains a list of the recently terminated contracts
\hyperref[appendix-c]{Appendix C}.) Its role is to take any specialized
action that is required before the contract is deleted (e.g., storing
contract slot information in the local database). No default action is
taken.

Information Acceptance Procedure: Called when an information message is
received. Its role is to determine what to do with the contents of the
information specification slot. The default is to store the information
in the local database.

Execution Procedure: Called when processing is started on the contract.
Its role is to perform the computation required to execute the task.

\section{ACKNOWLEDGMENT}\label{acknowledgment}

Some aspects of this work are being pursued in collaboration with R.
Davis at the Artificial Intelligence Laboratory of the Massachusetts
Institute of Technology, Cambridge. The contributions of B. Buchanan and
G. Wiederhold are also gratefully acknowledged.

\section{REFERENCES}\label{references}

{[}1{]} P. Brinch Hansen, Operating System Principles. Englewood Cliffs,
NJ: Prentice-Hall, 1973.\\
{[}2{]} D. J. Farber and K. C. Larson, ``The structure of the
distributed computing system---Software,'' in Proc. Symp. on
Comput.-Commun. Networks and Teletraffic, J. Fox, Ed. Brooklyn, NY:
Polytechnic Press, Polytechnic Inst. of Brooklyn, Apr.~1972,
pp.~539-545.\\
{[}3{]} D. J. Farber, J. Feldman, F. R. Heinrich, M. D. Hopwood, K. C.
Larson, C. Loomis, and L. A. Rowe, ``The distributed computing system,''
IEEE COMPCON Spring, 1973, pp.~31-34.\\
{[}4{]} C. Hewitt, ``Description and theoretical analysis (using
schemata) of PLANNER: A language for proving theorems and manipulating
models in a robot,'' Mass. Inst. Technol., Cambridge, MA, AI TR 258,
Apr.~1972.\\
{[}5{]} V. R. Lesser, ``Cooperative distributed processing,'' Dep.
Comput. Inform. Sci., Univ. of Massachusetts, Amherst, MA, COINS TR
78-7, May 1978.\\
{[}6{]} E.D. Sacerdoti, ``A structure for plans and behavior,'' SRI
Int., Menlo Park, CA, AIC TN 109, Aug.~1975.\\
{[}7{]} R. G. Smith and R. Davis, ``Applications of the contract net
framework: Distributed sensing,'' in Proc. ARPA Distributed Sensor Net
Symp., Pittsburgh, PA, Dec.~1978, pp.~12-20.\\
{[}8{]} R. G. Smith, ``A framework for problem solving in a distributed
processing environment,'' Dep. Comput. Sci., Stanford, CA,
STAN-CS-78-700 (HPP-78-28), Dec.~1978.\\
{[}9{]} R. G. Smith and R. Davis, ``Cooperation in distributed problem
solving,'' in Proc. /979 Int. Conf Cybern. Soc., Oct.~1979,
pp.~366-371.\\
{[}10{]} R. G. Smith, ``Applications of the contract net framework:
Search,'' in Proc. 1980 Nat. Conf. Canadian Soc. for Computational
Studies of Intell., May 1980, pp.~232-239.\\
{[}11{]} R. F. Sproul and D. Cohen, ``High-level protocols,'' Proc.
IEEE, vol.~66, pp.~1371-1386, Nov.~1978.\\
{[}12{]} W. Teitelman, INTERLISP Reference Manual. Palo Alto, CA: Xerox
Palo Alto Research Center, Dec.~1975.

\textit{[Image omitted: images/516170a0d45b42ee05dc77068fccc0e4717d6580103330bde5481d027a68338a.jpg]}

natural\_image

Portrait of a man with beard and short hair (no text or symbols visible)

Reid G. Smith (S'67-M'69-S'75-M'78) was born in Toronto, Ont., Canada,
on October 4, 1946. He received the B.Eng. and M. Eng. degrees in
electrical engineering from Carleton University, Ottawa, Ont., Canada in
1968 and 1969, respectively, and the Ph.D.~degree in electrical
engineering from Stanford University, Stanford, CA in 1979.

He has been associated with the Defence Research Establishment Atlantic
in Dartmouth, N.S., Canada since 1969. He is currently engaged in
research in artificial intelligence and distributed problem solving.

\end{document}
